\documentclass[11pt,a4paper]{article}

\usepackage{amsmath,amssymb,amsthm}
\usepackage{hyperref}
\usepackage{booktabs}
\usepackage{microtype}

\newtheorem{theorem}{Theorem}[section]
\newtheorem{proposition}[theorem]{Proposition}
\newtheorem{corollary}[theorem]{Corollary}
\newtheorem{lemma}[theorem]{Lemma}
\theoremstyle{definition}
\newtheorem{definition}[theorem]{Definition}
\newtheorem{remark}[theorem]{Remark}

\newcommand{\bstar}{\beta^*}
\newcommand{\dhat}{\hat{\delta}}
\newcommand{\Heis}{\mathrm{Heis}_3}
\newcommand{\ZqZ}{\mathbb{Z}/q\mathbb{Z}}
\newcommand{\Fq}{(\mathbb{Z}/q\mathbb{Z})^*}

\title{Fibre-Level Admissibility from Conjugate Weil Blocks:\\
Structural Derivation of the Pair Observable\\
\large O17 --- Spectral Admissibility Subprogram}

\author{J.~Beau\\
\small Independent Researcher}

\date{2026}

\begin{document}

  \maketitle

  \begin{abstract}
    O16~\cite{Beau2026a20} established that the physically relevant observable
    for the Weil-block capacity problem is defined on conjugate pairs
    $\{c, q-c\}$ rather than on individual spectral blocks,
    yielding the doubling $\delta_{\mathrm{pair}} = 2\,\delta_c \approx 7.44$
    consistent with the phenomenological range $[7.4, 10.6]$.
    However, the identification of conjugate pairs with the fibres of the
    non-injective projection~$\Pi$ was stated there as a structurally
    motivated hypothesis rather than a derived result.
    The present paper provides the structural derivation.
    We show that the raw Gram-Schmidt redundancy counts are
    \emph{exactly identical} when the initial supports coincide,
    and are structurally equivalent in general,
    establishing that conjugate blocks carry the same dynamical information.
    The proportionality constant $r(c,q)$ appearing in O16 is shown to
    be a normalisation artefact of the pipeline, not an intrinsic
    invariant of the representation.
    The only structurally robust quantity is therefore the equality
    $\delta_c = \delta_{q-c}$, which follows from the conjugation
    identity $\rho_{q-c} = \overline{\rho_c}$ and the norm-invariance
    of Gram-Schmidt orthogonalisation.
    This closes the logical gap left open in O16 and requalifies the
    O12--O15 results as correct analyses of a block-level observable
    that is not the appropriate physical unit.
    A key outcome of this work is the elimination of a misleading
    structural hypothesis: the factor $r(c,q)$, initially suspected
    in O16 to carry arithmetic significance (quadratic Gauss sums or
    character multiplicities), is shown here to be purely
    normalisation-dependent and gauge-like.
    This simplifies the theoretical structure by removing an
    apparent invariant that was not one.
  \end{abstract}

  \tableofcontents

  \section{The open problem after O16}
    \label{sec:problem}

    O16~\cite{Beau2026a20} identified the following situation.
    The Weil-block pipeline of O12--O15~%
    \cite{Beau2026a16,Beau2026a17,Beau2026a18,Beau2026a19}
    measures the redundancy observable $\sigma_c(n)$ for individual
    blocks indexed by a central character~$c \in \Fq$.
    The fitted exponent $\dhat_{\mathrm{exact}} \approx 3.72 \pm 0.06$
    is structurally below the phenomenological target $[7.4, 10.6]$,
    and no reweighting of block data can raise it above this value
    (O15 Corollary~4.8~\cite{Beau2026a19}).

    O16 proposed that the correct observable is defined on
    \emph{conjugate pairs} $\{c, q-c\}$:
    \begin{equation}
      \sigma_{\mathrm{pair}}(n) := \sigma_c(n)\,\sigma_{q-c}(n),
      \label{eq:sigma_pair}
    \end{equation}
    yielding $\delta_{\mathrm{pair}} = 2\,\delta_c \approx 7.44$.
    This was supported by:
    \begin{enumerate}
      \item the exact conjugation identity $\rho_{q-c} = \overline{\rho_c}$,
      \item numerical verification that $\delta_{q-c} = \delta_c$ to
      machine precision,
      \item the observation that $\sigma_{\mathrm{pair}}(n)$ follows a
      significantly cleaner power law than $\bar\sigma(n)$.
    \end{enumerate}

    However, O16 explicitly acknowledged (Remark~6.2) that the
    identification of conjugate pairs with the \emph{fibres of~$\Pi$}
    was not yet derived from first principles.
    Two questions were left open:

    \begin{itemize}
      \item[(Q1)] Why is the raw Gram-Schmidt dynamics identical for
      $\rho_c$ and $\rho_{q-c}$?
      Is this a consequence of $\rho_{q-c} = \overline{\rho_c}$,
      or does it require additional structure?

      \item[(Q2)] What is the structural origin of the
      constant~$r(c,q)$ satisfying
      $\sigma_{q-c}(n) = r(c,q)\,\sigma_c(n)$ for all $n \geq 1$?
      Is it an arithmetic invariant of $(c,q)$, or a pipeline artefact?
    \end{itemize}

    The present paper answers both questions.
    O17 is not a correction of O16, but the first derivation of the
    observable class dictated by~$\Pi$.

  \section{Identity of raw Gram-Schmidt dynamics}
    \label{sec:identity}

    \subsection{Setup}

      Let $q$ be an odd prime.
      The Weil representation $\rho_c$ of $\Heis(\ZqZ)$ at central
      character~$c \in \Fq$ acts on $\mathbb{C}^q$ by:
      \begin{equation}
        [\rho_c(a,b,t)\,\psi](x) \;=\; \omega^{c(bx + t)}\,\psi(x-a),
        \qquad \omega = e^{2\pi i/q}.
        \label{eq:weil_def}
      \end{equation}
      A block in the O12/O13 pipeline is specified by a triple
      $(c, b_1, b_2) \in \Fq \times (\ZqZ)^2$, with initial fingerprint
      vector:
      \begin{equation}
        v_0(c, b_1, b_2) \;:=\; \rho_c(b_1, b_2, 0)\,e_0
        \;=\; \omega^{c\,b_2 b_1}\,e_{b_1},
        \label{eq:v0}
      \end{equation}
      where $e_k$ denotes the $k$-th standard basis vector of $\mathbb{C}^q$.

      The Gram-Schmidt redundancy count at BFS shell~$n$ for this block is:
      \begin{equation}
        \tilde\sigma_c(n;\, b_1, b_2)
        \;:=\; \bigl|\bigl\{g \in \mathrm{Shell}_n :
        \rho_c(g)\,v_0(c,b_1,b_2)
        \in \mathrm{span}\bigl\{\rho_c(h)\,v_0 : h \in \mathrm{Shell}_{<n}\bigr\}\bigr\}\bigr|.
        \label{eq:raw_sigma}
      \end{equation}
      This is the \emph{raw} count: the number of Weil images at shell~$n$
      that are already in the Gram-Schmidt span accumulated over
      shells~$0, \ldots, n-1$.

    \subsection{Main results}

      The analysis splits into two levels: an exact equality under a support
      condition, and a structural isomorphism for arbitrary block parameters.

      \begin{theorem}[Identity of raw dynamics, restricted form]
        \label{thm:identity}
        Let $q$ be an odd prime and $c \in \Fq$.
        Suppose the two blocks share the same initial support, i.e.\ $b_1 = b_1'$.
        Then:
        \begin{equation}
          \tilde\sigma_{q-c}(n;\, b_1, b_2')
          \;=\; \tilde\sigma_c(n;\, b_1, b_2)
          \qquad \text{for all } n \geq 0.
          \label{eq:raw_identity}
        \end{equation}
      \end{theorem}

      \begin{proof}
        From equation~\eqref{eq:v0}, both initial vectors have support $\{b_1\}$:
        \[
          v_0(c,b_1,b_2) = \omega^{c\,b_2 b_1}\,e_{b_1},
          \qquad
          v_0(q-c,b_1,b_2') = \omega^{(q-c)\,b_2' b_1}\,e_{b_1}.
        \]
        These differ only by an overall unimodular scalar.
        The conjugation identity $\rho_{q-c} = \overline{\rho_c}$ gives,
        for every $g \in \Heis(\ZqZ)$:
        \[
          \rho_{q-c}(g)\,v_0(q-c,b_1,b_2')
          \;=\; \lambda\,\overline{\rho_c(g)\,v_0(c,b_1,b_2)},
        \]
        where $\lambda \in \mathbb{C}$ is a phase independent of~$g$.
        Since Gram-Schmidt orthogonalisation depends only on the
        linear independence structure of vectors, and since complex
        conjugation preserves this structure
        ($|\langle \bar u, \bar v \rangle| = |\langle u, v \rangle|$),
        the redundancy count of
        $\{\rho_{q-c}(g)\,v_0 : g \in \mathrm{Shell}_{\leq n}\}$
        equals that of
        $\{\rho_c(g)\,v_0 : g \in \mathrm{Shell}_{\leq n}\}$
        at every shell.
        Moreover, complex conjugation preserves linear independence relations
        over $\mathbb{C}$, ensuring that the span structure generated by the
        group orbit is unchanged: a set of vectors is linearly independent
        if and only if its complex conjugate is.
      \end{proof}

      \begin{remark}[Scope of Theorem~\ref{thm:identity}]
        \label{rem:scope}
        This equality holds at the level of raw redundancy counts under
        identical initial supports.
        In the general case of differing supports, the two dynamics remain
        isomorphic and yield identical asymptotic exponents, but the
        pointwise equality of counts may fail.
      \end{remark}

      \begin{corollary}[Structural equivalence for arbitrary blocks]
        \label{cor:structural_equiv}
        For arbitrary block parameters $(b_1, b_2)$ and $(b_1', b_2')$
        with $b_1 \neq b_1'$, the Gram-Schmidt dynamics of $\rho_c$ and
        $\rho_{q-c}$ are isomorphic: they produce the same sequence of
        redundancy counts up to a permutation of the generator ordering.
        In particular, the asymptotic exponents satisfy:
        \[
          \delta_{q-c} = \delta_c.
        \]
      \end{corollary}

      \begin{proof}
        The Gram-Schmidt process on $\{\rho_c(g)\,v_0\}$ depends on
        the representation $\rho_c$ and the initial vector $v_0$, but not
        on the specific labelling of~$e_{b_1}$.
        Under $\rho_{q-c} = \overline{\rho_c}$, the representation matrices
        are complex conjugates; the linear independence structure of any
        set of vectors is invariant under this transformation
        (since rank is preserved by conjugation).
        The resulting redundancy sequences may differ shell-by-shell by
        a permutation of generators within each shell, but they share the
        same asymptotic power-law exponent.
      \end{proof}

      \begin{remark}[What is structurally robust]
Theorem~\ref{thm:identity} establishes point-wise equality of
$\tilde\sigma_c$ and $\tilde\sigma_{q-c}$ under the support condition.
Corollary~\ref{cor:structural_equiv} establishes exponent equality
for arbitrary blocks.
The physically relevant quantity is the exponent, not the
shell-by-shell count: the unique structurally robust consequence
is $\delta_{q-c} = \delta_c$.
      \end{remark}

      \begin{remark}[Numerical confirmation]
For the pair $(c=1, c=60)$ at $q=61$ with blocks
$(b_1=27, b_2=2)$ and $(b_1=3, b_2=30)$, direct computation
from first principles gives identical raw redundancy sequences:
\[
  (\tilde\sigma_1(n),\, \tilde\sigma_{60}(n)) \;=\;
  (0, 2, 10, 34, 80, 162, 292, 474, 722, \ldots)
\]
for $n = 0, 1, 2, \ldots, 8$, confirming Theorem~\ref{thm:identity}
to numerical precision.
      \end{remark}

  \section{The factor $r(c,q)$ is a normalisation artefact}
    \label{sec:r_artefact}

    \subsection{Pipeline normalisation}

      The pipeline observable $\sigma_c(n)$ reported in O12--O15 is
      \emph{not} the raw count $\tilde\sigma_c(n)$ of
      equation~\eqref{eq:raw_sigma}.
      It is a normalised quantity:
      \begin{equation}
        \sigma_c(n;\, b_1, b_2)
        \;=\; \frac{\tilde\sigma_c(n;\, b_1, b_2)}{D(c, b_1, b_2, n)},
        \label{eq:sigma_normalised}
      \end{equation}
      where the denominator $D(c, b_1, b_2, n)$ depends on the block
      parameters and the shell index in a way determined by the
      pipeline implementation.

    \subsection{Consequence for $r$}

      Since $\tilde\sigma_{q-c}(n) = \tilde\sigma_c(n)$ by
      Theorem~\ref{thm:identity}, the ratio
      \begin{equation}
        r(c,q;\, b_1, b_2, b_1', b_2')
        \;=\; \frac{\sigma_{q-c}(n;\, b_1', b_2')}{\sigma_c(n;\, b_1, b_2)}
        \;=\; \frac{D(c, b_1, b_2, n)}{D(q-c, b_1', b_2', n)}
        \label{eq:r_ratio}
      \end{equation}
      is a \emph{ratio of normalisation denominators}.
      It is constant in~$n$ (as observed empirically) because the
      relative normalisation of two blocks is shell-independent
      once the block parameters are fixed.

      \begin{proposition}[$r$ depends on block parameters, not on $(c,q)$]
        \label{prop:r_artefact}
        The factor $r$ in equation~\eqref{eq:r_ratio} depends on the
        specific block instances $(b_1, b_2)$ and $(b_1', b_2')$,
        not solely on the character pair $(c, q-c)$.
        In particular, two blocks with the same character~$c$ but different
        $(b_1, b_2)$ parameters yield different values of~$r$ when paired
        with the same dual block.
      \end{proposition}

      \begin{proof}
        Direct computation for $q = 29$: the character $c = 11$ appears
        in the pipeline with instances $(b_1=18, b_2=19)$ and $(b_1=25, b_2=5)$.
        Paired with the block $c = 18$, $(b_1=6, b_2=15)$:
        \begin{align*}
          \sigma_{18}(n;\,6,15) / \sigma_{11}(n;\,18,19) &= 2 \quad \forall n \geq 1,\\
          \sigma_{18}(n;\,6,15) / \sigma_{11}(n;\,25,5)  &= 1 \quad \forall n \geq 1.
        \end{align*}
        The same dual block gives $r = 2$ or $r = 1$ depending on
        which instance of the primal block is used.
        Therefore $r$ is not a function of $(c, q)$ alone.
      \end{proof}

    \subsection{Consequences}

      Proposition~\ref{prop:r_artefact} has three immediate consequences.

      \begin{corollary}[Inertness of $r$ for the exponent]
        \label{cor:r_inert}
        Since $r$ is shell-independent and positive,
        $\log\sigma_{\mathrm{pair}}(n) = \log r + \log\sigma_c(n) + \log\sigma_{q-c}(n)$,
        so:
        \[
          \delta_{\mathrm{pair}}
          := -\frac{\mathrm{d}\log\sigma_{\mathrm{pair}}}{\mathrm{d}\log n}
          = \delta_c + \delta_{q-c} = 2\,\delta_c.
        \]
        The value of $r$ does not affect the exponent.
      \end{corollary}

      \begin{remark}[$r$ is gauge-like]
        \label{rem:r_gauge}
        The factor $r(c,q)$ is not only non-intrinsic to the representation,
        but is gauge-like in the following precise sense:
        it depends on the normalisation convention used to define $\sigma$,
        while leaving all asymptotic exponents invariant.
        Different choices of the denominator $D(c,b_1,b_2,n)$ yield
        different values of $r$, but the same $\delta_{\mathrm{pair}}$.
        An observable that is invariant under changes of normalisation
        convention is precisely what is needed for physical interpretation:
        the exponent $\delta_{\mathrm{pair}}$ is such an observable; $r$ is not.
        The factor $r(c,q)$ is therefore not an intrinsic invariant but a
        normalisation-dependent quantity, analogous to a gauge choice.
      \end{remark}

      \begin{corollary}[Closure of O16 Remark~6.4]
        \label{cor:remark_closure}
        The suggestion in O16 Remark~6.4 that $r(c,q)$ might originate
        from quadratic Gauss sums or character multiplicities is not
        supported.
        $r$ is determined by the pipeline normalisation~$D(c, b_1, b_2, n)$,
        which is a property of the measurement protocol, not of the
        representation.
      \end{corollary}

      \begin{corollary}[The structural invariant]
        \label{cor:invariant}
        The only structurally robust quantity is the equality
        \[
          \delta_{q-c} = \delta_c,
        \]
        which holds exactly as a consequence of $\rho_{q-c} = \overline{\rho_c}$
        and the norm-invariance of Gram-Schmidt orthogonalisation.
        This equality is independent of the pipeline normalisation.
      \end{corollary}

  \section{Why the pair is the correct observable}
    \label{sec:why_pair}

    \subsection{The hierarchy of observables}

      The O-series has progressively refined the observable:

      \begin{itemize}
        \item O12--O13: $\bar\sigma(n)$ = arithmetic mean over individual blocks.
        \emph{Observable class: block-level mean.}
        \item O14: corrected $\bar\sigma$ for normalisation and central-phase
        contributions. \emph{Observable class: corrected block-level mean.}
        \item O15: $R_n^{\mathrm{eff}}$ defined via exact block-level dynamics.
        \emph{Observable class: exact block-level.}
        \item O16--O17: $R_n^{\mathrm{pair}}$ = mean over conjugate pairs.
        \emph{Observable class: fibre-level.}
      \end{itemize}

      Each step corrects the previous one without invalidating it:
      O14--O15 correctly analyse the block-level observable;
      their results are exact within that observable class.
      The move to O16--O17 is not a correction of O14--O15 but
      a change of observable class.
      The observable was not mis-measured --- it was mis-identified.

    \subsection{Why fibre-level is the correct class}

      The projection $\Pi : \chi \to \mathcal{O}$ is non-injective by
      the Born--Infeld constraint~\cite{Beau2026a2}.
      An admissibility condition on $\mathcal{O}$ therefore corresponds
      to a condition on \emph{equivalence classes under $\Pi$}, not
      on individual preimage configurations.
      This is the key structural point: the non-injectivity of $\Pi$
      acts at the level of fibres $\Pi^{-1}(y)$, not at the level of
      individual blocks $\rho_c$.
      Observable quantities must therefore be defined on these
      equivalence classes.

      In the Weil realisation of $\Heis(\ZqZ)$, the minimal non-injectivity
      is the parity involution realised as $c \leftrightarrow q-c$.
      An admissible observable at character~$c$ must be jointly
      admissible with its conjugate at $q-c$: the two preimages
      $\rho_c$ and $\rho_{q-c}$ represent the same observable class,
      so admissibility must be assessed simultaneously for both.
      This identification is supported by
      Theorem~\ref{thm:identity} and Corollary~\ref{cor:structural_equiv}:
      conjugate representations $\rho_c$ and $\rho_{q-c}$ produce
      identical Gram-Schmidt dynamics, indicating that they correspond to
      indistinguishable observable content under $\Pi$.
      Two configurations with indistinguishable observable content
      belong to the same fibre of $\Pi$ by definition.

      Theorem~\ref{thm:identity} and Corollary~\ref{cor:structural_equiv}
      therefore show that the two preimages carry
      \emph{identical dynamical information} (same Gram-Schmidt dynamics).
      Measuring only one of them therefore captures exactly
      half of the dynamically accessible information about the fibre.
      The pair observable $\sigma_{\mathrm{pair}}(n) = \sigma_c(n) \cdot \sigma_{q-c}(n)$
      is the natural joint measure of fibre admissibility, corresponding
      to the product of individual admissibilities under the assumption
      of fibre independence.

    \subsection{Requalification of O14--O15}

      O14~\cite{Beau2026a18} and O15~\cite{Beau2026a19} are not
      erroneous.
      They correctly characterise what happens when admissibility is
      assessed at the block level rather than the fibre level.
      The aggregation no-go (O15 Corollary~4.8) establishes that
      no reweighting of block data recovers the fibre-level exponent:
      this is precisely because block-level observables are
      half-fibres, and no averaging over half-fibres reconstructs
      the full-fibre observable.

      \begin{remark}[Complementarity of O15 and O17]
O15 establishes the impossibility of recovering the correct
exponent from block-level data.
O17 establishes what the correct exponent is and why it requires
fibre-level data.
The two results are complementary: O15 closes the block-level
programme; O17 opens the fibre-level programme.
      \end{remark}

  \section{Summary and open directions}
    \label{sec:summary}

    \subsection{What is established}

      \begin{enumerate}
        \item The raw Gram-Schmidt redundancy dynamics are exactly identical for
        $\rho_c$ and $\rho_{q-c}$ under equal initial support
        (Theorem~\ref{thm:identity}), and structurally isomorphic for
        arbitrary block parameters (Corollary~\ref{cor:structural_equiv}),
        as a consequence of $\rho_{q-c} = \overline{\rho_c}$ and
        norm-invariance of the Gram-Schmidt process.
        The unique structurally robust consequence is $\delta_{q-c} = \delta_c$.

        \item The factor $r(c,q)$ in O16 is a pipeline normalisation
        artefact (Proposition~\ref{prop:r_artefact}), not an
        arithmetic invariant of the representation.
        It is gauge-like (Remark~\ref{rem:r_gauge}): it depends on the
        normalisation convention, while leaving all asymptotic
        exponents invariant.
        The suggestion in O16 Remark~6.4 that $r$ might originate from
        quadratic Gauss sums or character multiplicities is not supported.

        \item The equality $\delta_{q-c} = \delta_c$ is the unique
        structurally robust consequence, independent of pipeline
        normalisation (Corollary~\ref{cor:invariant} and
        Corollary~\ref{cor:structural_equiv}).

        \item The pair observable is the correct unit of admissibility
        because the non-injective projection $\Pi$ requires fibre-level
        assessment, and conjugate pairs are the fibres of~$\Pi$
        in the Weil realisation.

        \item O14--O15 are correct analyses of the block-level observable.
        The move to fibre-level is a change of observable class,
        not a correction of previous errors.
      \end{enumerate}

    \subsection{Open directions}

      \paragraph{Primary open problem: derivation of the conjugate pairing.}
        The identification of $c \leftrightarrow q-c$ as the minimal
        non-injectivity of $\Pi$ remains a structural hypothesis
        (O16 Remark~6.2, O17 Section~\ref{sec:why_pair}).
        A formal derivation within the Cosmochrony framework ---
        showing that the parity involution $c \mapsto q-c$ is a
        \emph{necessary} consequence of the projection structure~$\Pi$
        and the Born--Infeld constraint --- would promote the main result
        from a well-supported structural result to a theorem.
        This is the primary open problem of the fibre-level programme.

      \paragraph{Secondary: canonical normalisation of $\sigma$.}
        The pipeline normalisation $D(c, b_1, b_2, n)$ is currently
        implicit in the O12/O13 implementation.
        Making it explicit and defining a canonical normalisation
        independent of $(b_1, b_2)$ would eliminate $r$ entirely from
        the pair observable, producing a version of $\sigma_{\mathrm{pair}}$
        that is genuinely independent of the measurement convention.
        This is a technical rather than conceptual question.

      \paragraph{Upper bound of the phenomenological window.}
        Numerical exploration (post-O16) shows that the spectrum
        of $\delta_{\mathrm{pair}}$ has a hard lower bound near~$7.4$
        and extends significantly above~$10.6$.
        The mechanism selecting $[7.4, 10.6]$ as the physical
        sub-spectrum is not yet identified; it appears to be a
        dynamical persistence criterion rather than a spectral one.

  \section*{Acknowledgements}

  The author thanks the extended Cosmochrony programme discussions
  that led to the identification of the normalisation origin of~$r$
  and the structural derivation of the fibre-level observable.

  \bibliographystyle{plain}
  \begin{thebibliography}{99}
    \input{cosmochrony-bibliography}
  \end{thebibliography}

\end{document}
